\section{不变函数与等变映射：对称性怎样进入计算}
\label{sec:04-Discrete-Mathematics/06-Groups-and-Symmetry/04}

一个预测分子能量的模型跑通了。你把同一个分子的原子在输入文件里换个书写顺序，再跑一遍，预测值挪了 \(0.3\) 电子伏。分子一个原子都没动，动的只是编号。

另一个场景更常见：手写数字识别器在测试集上很稳，把每张图整体右移三个像素之后，准确率掉了十几个百分点。数字还是那个数字。

这两次失败长得一样。输入沿着某条轨道移动了一下，输出本不该跟着动，却动了。上一节已经把``沿轨道移动''说清楚了，现在缺的是另一半：把``输出应该怎样响应''写成对函数的约束，写进模型，而不是指望它自己学会。

\subsection{两种响应方式}

设群 \(G\) 作用在输入空间 \(X\) 上。函数 \(f:X\to Y\) 称为 \(G\)-不变（invariant），若

\[
f(g\cdot x)=f(x)\qquad\text{对所有 }g\in G,\ x\in X.
\]

等价说法是 \(f\) 在每条轨道上取常值——它看不见轨道内部的差别。向量长度对正交群不变，\(\norm{Qx}=\norm{x}\)；一张照片的类别标签对旋转不变；一个分子的总能量对原子编号的置换不变。

若输出空间 \(Y\) 上也有 \(G\) 的作用，映射 \(F:X\to Y\) 称为 \(G\)-等变（equivariant），若

\[
F(g\cdot x)=g\cdot F(x)\qquad\text{对所有 }g\in G,\ x\in X.
\]

输入被对称操作怎样改动，输出就照样改动一遍。

\begin{center}
\begin{tikzpicture}[x=1cm,y=1cm,font=\small,>=stealth]
  \node (a) at (0,1.6) {\(x\)};
  \node (b) at (3.4,1.6) {\(g\cdot x\)};
  \node (c) at (0,0) {\(F(x)\)};
  \node (d) at (3.4,0) {\(g\cdot F(x)\)};
  \draw[->] (0.35,1.6) -- (2.85,1.6) node[midway,above] {\(g\)};
  \draw[->] (0,1.3) -- (0,0.35) node[midway,left] {\(F\)};
  \draw[->] (3.4,1.3) -- (3.4,0.35) node[midway,right] {\(F\)};
  \draw[->] (0.75,0) -- (2.5,0) node[midway,above] {\(g\)};
  \node[right,black!70] at (4.6,0.8) {两条路径到达同一处：\(F\) 等变};
  \draw[black!35] (-0.9,-0.6) rectangle (4.3,2.2);
\end{tikzpicture}

\emph{等变说的是这张方图可交换：先作用后映射，等于先映射后作用。不变是它的特例——让 \(G\) 在 \(Y\) 上平凡作用（\(g\cdot y=y\)），下面那条横箭头就退化成恒等，条件变成 \(F(g\cdot x)=F(x)\)。}
\end{center}

区分这两者不是术语洁癖，它直接决定网络该怎么搭。把一张图像旋转 \(90^\circ\)，分类标签不该变（要不变），但语义分割的掩码必须跟着转同样的角度（要等变）。要求错了级别，模型会在对称方向上系统性地出错，而且这种错误在同分布的测试集上往往看不出来。

\subsection{卷积等变，池化不变}

深度学习里最熟的一对例子恰好把两个概念摆在同一条流水线上。

\begin{center}
\begin{tikzpicture}[x=1cm,y=1cm,font=\footnotesize]
  \node[black!75] at (2.15,5.15) {原输入};
  \node[black!75] at (8.05,5.15) {整体右移两格};
  \foreach \bx in {0,5.9} {
    \begin{scope}[shift={(\bx,0)}]
      \draw[black!45] (-0.1,3.6) -- (4.5,3.6);
      \draw[black!45] (-0.1,1.8) -- (4.5,1.8);
      \draw[black!45] (-0.1,0) -- (4.5,0);
    \end{scope}
  }
  \foreach \i/\h in {0/0.08,1/0.14,2/0.85,3/1.15,4/0.55,5/0.14,6/0.08,7/0.08}
    \fill[blue!40] (0.55*\i,3.6) rectangle +(0.42,\h);
  \foreach \i/\h in {0/0.10,1/0.45,2/0.95,3/1.05,4/0.45,5/0.12,6/0.08,7/0.08}
    \fill[green!45] (0.55*\i,1.8) rectangle +(0.42,\h);
  \fill[orange!55] (1.65,0) rectangle +(0.42,1.05);
  \begin{scope}[shift={(5.9,0)}]
    \foreach \i/\h in {0/0.08,1/0.08,2/0.08,3/0.14,4/0.85,5/1.15,6/0.55,7/0.14}
      \fill[blue!40] (0.55*\i,3.6) rectangle +(0.42,\h);
    \foreach \i/\h in {0/0.08,1/0.08,2/0.10,3/0.45,4/0.95,5/1.05,6/0.45,7/0.12}
      \fill[green!45] (0.55*\i,1.8) rectangle +(0.42,\h);
    \fill[orange!55] (1.65,0) rectangle +(0.42,1.05);
  \end{scope}
  \draw[black!55,dashed] (1.86,3.5) -- (1.86,4.95);
  \draw[black!55,dashed] (1.86,1.7) -- (1.86,3.05);
  \draw[black!55,dashed] (7.96,3.5) -- (7.96,4.95);
  \draw[black!55,dashed] (7.96,1.7) -- (7.96,3.05);
  \node[left,black!80] at (-0.25,4.0) {输入};
  \node[left,black!80] at (-0.25,2.2) {卷积特征};
  \node[left,black!80] at (-0.25,0.5) {全局池化};
\end{tikzpicture}

\emph{同一段一维信号，右边是把它整体右移两格的版本。卷积特征的峰跟着挪了两格（虚线标出峰位），这是平移等变；全局池化只报最大响应的大小，两边给出同一根柱子，这是平移不变。位置信息在中间层一直保留着，直到最后一步才被主动丢掉。}
\end{center}

这张图解释了标准配方为什么长成那样：主干层保持等变，读出层做不变化。若从第一层起就不变，位置信息在网络的最前面就没了，后面的层再也拿不回来，检测和分割这类需要``东西在哪''的任务直接做不了。反过来，若整条链都等变而没有读出层，分类头就会随着输入平移而给出漂移的分数。\hyperref[sec:14-Deep-Learning/04-Convolutional-Networks/01]{``卷积把平移结构写进线性映射''}里说的``逐层保持位置、最后不随位置变化''，在这里有了群论的写法。

\subsection{等变约束就是参数共享}

把对称写进模型，落到实处只有一个机制：绑定参数。设群在输入与输出上的作用都由矩阵给出，记作 \(\rho_X(g)\) 与 \(\rho_Y(g)\)。线性层 \(y=Wx\) 等变，当且仅当

\[
W\rho_X(g)=\rho_Y(g)W\qquad\text{对所有 }g\in G.
\]

这是一组关于 \(W\) 的线性方程。解出来的解空间维数，就是这一层剩下的自由参数个数——对称性把参数量压缩了多少，可以直接算。

取 \(G\) 为长度 \(n\) 的循环平移群，\(\rho\) 是循环移位矩阵。与所有移位矩阵交换的矩阵恰好是循环矩阵：每一行是上一行右移一格。自由参数从 \(n^2\) 掉到 \(n\)，而循环矩阵乘向量就是循环卷积。卷积不是有人拍脑袋想出来的算子，它是上面那个方程对平移群的通解。

取 \(G=S_n\) 作用在 \(\R^n\) 上（置换坐标），\(\rho_X=\rho_Y\) 是置换矩阵。与所有置换矩阵交换的矩阵只有两个自由度：

\[
W=a\,I+b\,\mathbf{1}\mathbf{1}^{\trans}.
\]

\(n^2\) 个参数压成 2 个。这一层做的事情是``每个位置保留自己一份，再加上全体的和''，正是集合上的深度网络（Deep Sets）与消息传递图网络里最基本的那一步：邻居聚合用求和或平均，与邻居的排列顺序无关，因而对节点置换等变；最后的读出层再把所有节点求和，得到置换不变的图级输出。点云模型、分子模型的骨架都是这个形状。

群越大，方程组越紧，剩下的参数越少。这带来两重好处：待估参数少，需要的样本就少；模型在整条轨道上的行为被一次性钉死，不需要靠数据去覆盖。粗略地说，硬等变约束相当于让模型在商掉对称之后的 \(X/G\) 上学习，而 \(X/G\) 比 \(X\) 小 \(\abs{G}\) 倍左右。

数据增强追求的是同一件事，但走的是软路线：训练时对每个样本随机施加群元素，把同一条轨道上的多个点喂进去，指望模型自己归纳出不变。两条路的差别在保证的强度上。增强只在训练分布覆盖到的地方鼓励不变，覆盖不到的角落没有约束；等变结构在全部输入上恒等成立，包括没见过的那些。代价是它要求对称假设一开始就是对的。实践中两者常叠加：结构上写死能确定的那部分对称，剩下的近似对称交给增强。

\subsection{不变量够不够分辨轨道}

一个不变函数把整条轨道压成一个值。反过来问：一组不变量能不能把不同的轨道分开？这是不变量理论的基本问题，也是特征工程反复撞上的那堵墙。

造不变函数很容易。对有限群，取任意 \(f\) 做群平均

\[
\bar f(x)=\frac1{\abs{G}}\sum_{g\in G}f(g\cdot x)
\]

就得到一个不变函数。容易归容易，得到的东西可能平得什么也分不出——常函数也是完美的不变函数。

好的例子是 \(S_n\) 作用在 \(\R^n\) 上。基本对称多项式

\[
e_1=\sum_i x_i,\quad e_2=\sum_{i<j}x_ix_j,\quad\ldots,\quad e_n=x_1x_2\cdots x_n
\]

能分离轨道：两个向量的坐标只差一个排列，当且仅当所有 \(e_k\) 相等——因为这组数恰好是多项式 \(\prod_i(t-x_i)\) 的系数，而系数决定根的重集。``只看集合不看顺序''这个日常要求，被 \(n\) 个多项式完全参数化了。

不是每种情形都这么幸运。消息传递图网络对节点置换等变，读出层给出置换不变的图表示，可这个不变量分不开所有非同构的图：它的判别能力有一个上界，与一维 Weisfeiler--Lehman 颜色细化算法相同，存在成对的非同构图被映到完全相同的表示。换句话说，函数是不变的，但轨道没被分开，两张结构不同的图在模型眼里一模一样。\hyperref[sec:04-Discrete-Mathematics/03-Graph-Theory/10]{``图同构与可达性：结构不依赖画法''}问的正是两张图是否落在同一条轨道上；不变量给出的是必要条件，不是充分条件。诊断办法也顺着这条线走：模型在某类样本上总是给出雷同的输出时，先检查是不是它的不变量把本该分开的轨道合并了。

\subsection{对称假设选错的代价是永久的}

把对称写进结构，好处是保证，坏处也是保证——它在所有输入上都成立，包括你希望它别成立的地方。

旋转 \(180^\circ\) 把 \(6\) 变成 \(9\)。文字识别任务的对称群里不能放旋转，否则模型被结构性地禁止区分这两个类。医学影像的左右位置常常携带临床含义，强加镜像不变等于提前删掉可能救命的信息。航拍图里``上''是一个真实存在的方向，地平线总是接近水平；对这类数据强加旋转等变，会把这条免费的先验一并抹掉，模型反而更难学。

判断的顺序值得固定下来。先问哪个群在作用——平移、旋转、反射，还是置换；再问它作用在输入、输出还是中间表示上；再问任务要的是不变还是等变，标签该钉住还是该跟着动；最后问同一条轨道内的差异是否当真与任务无关。最后这一问最容易被跳过，也最容易出事：合并轨道等于宣布轨道里的差异没有信息，这是一条关于数据和任务的实质性断言，不是一个安全的默认值。

数据只是近似对称时，硬约束与软增强的取舍要重新算。约束下得越硬，收益（参数少、样本省）和风险（抹掉有用方向）同步放大。一个折中是把对称当成正则项而非结构：在损失里惩罚 \(\norm{F(g\cdot x)-g\cdot F(x)}\)，强度可调，对称偏离多少由数据说了算。

这一章走完的路径是：从一批``看不出区别''的动作出发得到群，让群作用到对象上得到轨道，用轨道--稳定子和 Burnside 数清轨道，再把轨道要求翻译成对函数的不变与等变约束。上面反复出现的 \(\rho(g)\) 还只是被当作记号使用——群怎样系统地变成矩阵、这些矩阵能拆到多细，是\hyperref[ch:035]{``群表示论：对称的线性结构''}的主题，等变网络的设计空间也在那里才能完整地写出来。
